\section{Билет 50. Адиабатический процесс. Уравнение адиабаты без вывода. Постоянная адиабаты. Работа газа в адиабатическом процессе с выводом.}

\begin{definition}
\textbf{Адиабатический процесс} протекает без теплообмена с окружающей средой ($Q = 0$). Связь между давлением $p$ и объёмом $V$ идеального газа в таком процессе описывается \textbf{уравнением Пуассона}:
\begin{equation}
    p V^\gamma = \text{const}, \label{eq:poisson_50}
\end{equation}
где $\gamma = C_p/C_V$ --- \textbf{постоянная адиабаты} (показатель адиабаты), зависящая от молекулярной структуры газа: $\gamma = (i+2)/i > 1$.
\end{definition}

\begin{theorem}[Работа газа в адиабатическом процессе]
Согласно первому началу термодинамики при $Q=0$, работа газа совершается исключительно за счёт убыли его внутренней энергии:
\begin{equation}
    A = -\Delta U = -(U_2 - U_1) = U_1 - U_2. \label{eq:work_energy_50}
\end{equation}
Внутренняя энергия идеального газа $U = \nu C_V T$. Следовательно:
\begin{equation}
    A = \nu C_V (T_1 - T_2). \label{eq:work_Cv_50}
\end{equation}
Выразим $C_V$ через $\gamma$ и $R$, используя соотношение Майера $C_p - C_V = R$ и определение $\gamma = C_p/C_V$:
\begin{equation}
    \gamma C_V - C_V = R \quad \Rightarrow \quad C_V = \frac{R}{\gamma - 1}.
\end{equation}
Подставляем в \eqref{eq:work_Cv_50}:
\begin{equation}
    A = \nu \frac{R}{\gamma - 1} (T_1 - T_2). \label{eq:work_T_50}
\end{equation}
Используя уравнение состояния $\nu R T = p V$, заменяем $\nu R T_1 = p_1 V_1$ и $\nu R T_2 = p_2 V_2$. Получаем формулу для работы через начальные и конечные макроскопические параметры:
\begin{equation}
    A = \frac{p_1 V_1 - p_2 V_2}{\gamma - 1}. \label{eq:work_pV_50}
\end{equation}
\end{theorem}

\begin{property}
Из формулы \eqref{eq:work_pV_50} следуют важные свойства:
\begin{enumerate}
    \item При адиабатном \textbf{расширении} ($V_2 > V_1$) газ совершает положительную работу ($A > 0$). При этом $p_2 V_2 < p_1 V_1$, следовательно $T_2 < T_1$ --- газ \textbf{охлаждается}.
    \item При адиабатном \textbf{сжатии} ($V_2 < V_1$) работа внешних сил над газом положительна, а работа газа отрицательна ($A < 0$). Температура газа \textbf{повышается} ($T_2 > T_1$).
    \item Используя уравнение Пуассона $p_1 V_1^\gamma = p_2 V_2^\gamma$, работу можно выразить через отношение объёмов:
    \begin{equation}
        A = \frac{p_1 V_1}{\gamma - 1} \left[ 1 - \left(\frac{V_1}{V_2}\right)^{\gamma - 1} \right].
    \end{equation}
\end{enumerate}
\end{property}

\begin{remark}
\textbf{Графическое определение работы.} На $pV$-диаграмме работа численно равна площади под кривой адиабаты. Поскольку адиабата круче изотермы, при расширении от одного объёма до другого адиабатная работа \textbf{меньше} изотермической (при тех же начальных параметрах), так как давление падает быстрее из-за охлаждения газа.
\par\noindent\textbf{Практические примеры:} охлаждение воздуха при выходе из баллона (образование тумана), резкий нагрев дизельного топлива при адиабатном сжатии в цилиндре двигателя, работа турбин и компрессоров в термодинамических циклах.
\end{remark}

\begin{figure}[htbp]
    \centering
    \begin{tikzpicture}[
        scale=1.2,
        >=Stealth,
        font=\small,
        thick
    ]
    
    \node[font=\bfseries] at (3, 4.6) {Адиабатический процесс ($Q = 0$)};
    
    % === ОСИ ===
    \draw[->] (0, 0) -- (5.5, 0) node[right, font=\normalsize] {$V$};
    \draw[->] (0, 0) -- (0, 4) node[above, font=\normalsize] {$p$};
    \node[below left] at (0, 0) {$O$};
    
    % === ПАРАМЕТРЫ ===
    \def\Vone{1.5}
    \def\Vtwo{4.0}
    \def\Pone{3.33}
    \def\gamma{1.4}
    \def\Ciso{5.0}           % p = 5/V (изотерма)
    \def\Cadia{5.75}         % p = 5.85/V^1.4 (адиабата)
    
    % === ИЗОТЕРМА (синяя, менее крутая) ===
    \draw[blue, thick, domain=1.2:5.0, samples=100] 
        plot (\x, {\Ciso/\x});
    
    % === АДИАБАТА (красная, более крутая) ===
    \draw[red, very thick, domain=1.2:5.0, samples=100] 
        plot (\x, {\Cadia/pow(\x, \gamma)});
    
    % === ЗАШТРИХОВАННЫЕ ПЛОЩАДИ ===
    % Площадь под изотермой (голубая)
    \fill[blue!15, opacity=0.5, domain=\Vone:\Vtwo, samples=50] 
        plot (\x, {\Ciso/\x}) -- (\Vtwo, 0) -- (\Vone, 0) -- cycle;
    
    % Площадь под адиабатой (красная, поверх)
    \fill[red!15, opacity=0.5, domain=\Vone:\Vtwo, samples=50] 
        plot (\x, {\Cadia/pow(\x, \gamma)}) -- (\Vtwo, 0) -- (\Vone, 0) -- cycle;
    
    % === ТОЧКА 1 (общая для обеих кривых) ===
    \fill[black] (\Vone, \Pone) circle (3pt);
    \node[above left] at (\Vone, \Pone) {$1$};
    
    % === ТОЧКА 2 (изотерма) ===
    \def\PtwoIso{1.25}  % 5/4.0
    \fill[blue] (\Vtwo, \PtwoIso) circle (3pt);
    
    % === ТОЧКА 2 (адиабата) ===
    \def\PtwoAdia{0.84}  % 5.85/4^1.4
    \fill[red] (\Vtwo, \PtwoAdia) circle (3pt);
    
    % === ВЕРТИКАЛЬНЫЕ ЛИНИИ ===
    \draw[dashed, gray] (\Vone, 0) -- (\Vone, \Pone);
    \draw[dashed, blue] (\Vtwo, 0) -- (\Vtwo, \PtwoIso);
    \draw[dashed, red] (\Vtwo, 0) -- (\Vtwo, \PtwoAdia);
    
    % === ПОДПИСИ ОСЕЙ ===
    \node[below] at (\Vone, 0) {$V_1$};
    \node[below] at (\Vtwo, 0) {$V_2$};
    \node[left] at (0, \Pone) {$p_1$};
    
    % === СТРЕЛКИ НАПРАВЛЕНИЯ ===
    \draw[->, blue, thick] (2.5, 2.0) -- (3.0, 1.67);
    \draw[->, red, thick] (2.5, 1.7) -- (3.0, 1.3);
    
    % === ПОДПИСИ КРИВЫХ ===
    \node[blue, font=\footnotesize] at (4.5, 1.5) {Изотерма};
    \node[red, font=\footnotesize] at (4.5, 0.5) {Адиабата};
    
    % === БЛОК СРАВНЕНИЯ РАБОТ ===
    \node[draw, fill=yellow!10, rounded corners, align=left, font=\small, inner sep=5pt] at (3, -1.2) 
        {Работа при расширении:\\[4pt]
         $A_{\text{изот}} > A_{\text{адиаб}}$\\[4pt]
         (адиабата круче изотермы)};
    
    % === БЛОК УРАВНЕНИЙ ===
    \node[draw, fill=green!10, rounded corners, align=left, font=\small, inner sep=4pt] at (4.3, 3.5) 
        {Изотерма: $pV = \text{const}$\\[4pt]
         Адиабата: $pV^\gamma = \text{const}$\\[4pt]
         $\gamma = \frac{C_p}{C_V} > 1$};
    
    \end{tikzpicture}
    \caption{Сравнение адиабатического и изотермического процессов на pV-диаграмме. Обе кривые начинаются в точке 1. Адиабата (красная) круче изотермы (синяя), поэтому при расширении от $V_1$ до $V_2$ работа газа в адиабатическом процессе меньше: $A_{\text{адиаб}} < A_{\text{изот}}$.}
    \label{fig:adiabatic_vs_isothermal}
\end{figure}

